\chapter{Inference: Methodology and Results}
\label{chap:main-inference}

The estimation results show that early stopping can protect slope recovery
from weakly identified directions.  Hypothesis testing places a different
demand on the fit: reducing the remaining empirical moment residual makes
the fitted statistic approximate a direct sampling-based statistic.  Following
\textcite{babii_carrasco_tsafack_2026}, this chapter studies how continued
moment-residual CG--FPLS iterations support inference, and assesses stopping,
null rejection, power, and confidence sets in the three simulation models.
These results concern CG--FPLS inference; they do not compare inferential
procedures for all four estimation methods.

\section{The hypothesis and the role of stopping}
\label{main:sec-inference-statistic}

For the centred model $Y^c=\langle X^c,\beta\rangle+\varepsilon$ with
$\mathbb E(\varepsilon\mid X)=0$, the identifiable hypothesis is
$H_0:K\beta=Kb$, where $b$ is a specified candidate slope.  It becomes the
literal slope hypothesis $H_0:\beta=b$ when both functions belong to
$\operatorname{Null}(K)^\perp$.  A difference in $\operatorname{Null}(K)$
cannot be detected.  Choosing $b=0$ tests for a functional effect; to measure
null rejection, the simulations use their known generating slope as $b$.

The hypothesis concerns the complete candidate function, rather than its
value at a particular wavelength or time point.  Rejection means that the
candidate is inconsistent with the functional moment equation at the chosen
significance level.  Non-rejection means that the evidence is insufficient
to rule it out; it does not establish that the candidate is correct.  This
distinction matters in an ill-posed problem, where substantial changes in
weak covariance directions can have little effect on the observed moment.

Using the empirical moments from Chapter~\ref{chap:main-methods}, define
\begin{equation}
 \begin{aligned}
 u_n(b)&=\widehat r-\widehat Kb,\qquad
 e_m=\widehat r-\widehat K\widehat\beta_m^{\mathrm{CG}},\\
 \mathcal T_{n,m}(b)&=n\|\widehat K(\widehat\beta_m^{\mathrm{CG}}-b)\|^2.
 \end{aligned}
 \label{main:eq-inference-statistic}
\end{equation}
Here $e_m$ is the remaining moment residual.  Applying $\widehat K$ avoids
requiring a root-$n$ limit for the unstable, unweighted slope estimate.
Since $\widehat K(\widehat\beta_m^{\mathrm{CG}}-b)=u_n(b)-e_m$, comparison
with the direct statistic $\mathcal S_n(b)=n\|u_n(b)\|^2$ gives
\begin{align}
 \mathcal T_{n,m}(b)-\mathcal S_n(b)
 &=-2n\langle u_n(b),e_m\rangle+n\|e_m\|^2,\label{main:eq-inference-difference}\\
 |\mathcal T_{n,m}(b)-\mathcal S_n(b)|
 &\leq2\sqrt n\|u_n(b)\|\sqrt n\|e_m\|+n\|e_m\|^2.\nonumber
\end{align}
Under the null, $\sqrt n\|u_n(b)\|=O_{\mathbb P}(1)$.  Consequently,
\begin{equation}
 \sqrt n\|e_{m_n}\|=o_{\mathbb P}(1)
 \quad\Longrightarrow\quad
 \mathcal T_{n,m_n}(b)=\mathcal S_n(b)+o_{\mathbb P}(1).
 \label{main:eq-inference-late}
\end{equation}
This is a sufficient condition for iteration error to disappear from the
limiting null distribution.  An estimation rule permitting a constant
multiple of $n^{-1/2}$ for $\|e_m\|$ does not ensure it.  Further iterations
can therefore improve the test approximation while worsening slope
estimation.  The direct statistic itself requires no CG fit; its role is
both an inferential construction and a benchmark for assessing the fitted
statistic.  No fixed component count guarantees
\eqref{main:eq-inference-late} across models and sample sizes.

The comparison with $\mathcal S_n$ also clarifies what continuing CG
accomplishes.  It connects an estimate obtained by regularised iteration
to a statistic whose fluctuation can be analysed directly.  It does not
make inverse estimation unnecessary when the objective is to recover the
slope, and it does not make the limiting test automatically usable without
a critical value.  Accurate solution of the empirical equation and accurate
calibration of the resulting statistic are separate requirements.  The
experiments below examine them separately so that one cannot conceal a
failure of the other.

\section{Null law and calibration}
\label{main:sec-inference-calibration}

Assume independent, identically distributed observations, $\mathbb E\|X^c\|^4<\infty$, and bounded
conditional error variance.  Under the null and known population means,
$u_n(b)=n^{-1}\sum_i\varepsilon_iX_i^c$.  Sample centring adds a term of
order $O_{\mathbb P}(n^{-1})$, which is negligible on the root-$n$ scale.
The Hilbert-space central limit theorem and
\eqref{main:eq-inference-late} give
\begin{equation}
 \sqrt n\,u_n(b)\rightsquigarrow G\sim\mathcal N_{\mathcal H}(0,V),
 \qquad
 \mathcal T_{n,m_n}(b)\rightsquigarrow
 \|G\|^2=\sum_j\omega_jZ_j^2,
 \label{main:eq-inference-null-law}
\end{equation}
where $V=\mathbb E\{\varepsilon^2(X^c\otimes X^c)\}$, the $\omega_j$ are
its positive eigenvalues, and the $Z_j$ are independent standard normals.
The moment conditions make $V$ trace class.  Under conditional
homoskedasticity, $V=\sigma^2K$ and $\omega_j=\sigma^2\lambda_j$;
under heteroskedasticity, $V$ and $K$ need not share eigenfunctions.
Rejecting above the $(1-\alpha)$ quantile of this non-degenerate limiting
law gives asymptotic size $\alpha$, provided the residual condition holds.
Proofs appear in Appendix~\ref{app:mathematical-proofs}.

The simulations retain the three models of
Chapter~\ref{chap:main-simulation}, with 100 Gaussian score terms, a
200-point grid, and error standard deviation one.  Each design point uses
5,000 samples.  Principal null-size results use $n=100$, $m=70$, and
critical values from 50,000 draws of
$\sum_{j=1}^{100}\sigma^2\lambda_jZ_j^2$.  This \emph{oracle} calibration
uses known population parameters, rather than estimating them from an
observed dataset.

The continuous-model oracle law also differs from the exact fixed-grid
limit.  If $\Phi_{jt}=\phi_j(s_t)$ contains the generating basis values and
$G_T=T^{-1}\Phi\Phi^\top$, the grid covariance eigenvalues are those of
$\Lambda^{1/2}G_T\Lambda^{1/2}$, with
$\Lambda=\operatorname{diag}(\lambda_1,\ldots,\lambda_{100})$.
The endpoint-inclusive grid has $G_T\neq I$.  Increasing $n$ at fixed
$T=200$ does not remove this discrepancy.  Departures from the oracle
reference can therefore reflect finite-sample effects, grid approximation,
and uncertainty in the simulated critical value.  The power study instead
uses empirical null quantiles, separately for each model and sample size.
Neither calibration is available directly from a single observed dataset.

For an application, the unknown score covariance must instead be estimated.
For a specified null candidate, one possible route uses residuals
$\widehat\varepsilon_i(b)=Y_i^c-\langle X_i^c,b\rangle$ and the operator
$\widehat V=n^{-1}\sum_i\widehat\varepsilon_i(b)^2
(X_i^c\otimes X_i^c)$, then simulates the weighted chi-square law using
its estimated eigenvalues.  Bootstrap calibration provides another route.
These possibilities motivate a feasible procedure, but their finite-sample
performance is not established by an experiment that supplies the true
population weights.  The present design deliberately makes that limitation
visible: its oracle calculation checks the theoretical approximation, while
its empirical-null power calculation isolates detection from the observed
oracle size distortion.

\section{Stopping and null rejection in the simulations}
\label{main:sec-inference-size}

The longer inference path uses two-pass reorthogonalisation of the Krylov
basis and a QR solve that minimises
$\|\widehat r-\widehat KQ_ma\|$.  In exact arithmetic this has the same
minimum-moment-residual objective as the released CG--FPLS recurrence.
It differs from Arnoldi FPLS, which minimises response error.  Numerical
stabilisation is particularly relevant when many iterations are taken.

A separate stopping experiment compares fixed dimensions, a maximum-rank
fit, and an adaptive estimation-oriented variant.  The variant retains
$\tau=1.01$ and $\delta_{\mathrm{stop}}=0.1$, but limits the variance pilot
to 20 components rather than the estimation study's 70 and uses the
stabilised path.  It is therefore not an exact rerun of the estimation
procedure.  Its selected median dimensions are one, two, and one in
Models~1--3.  All 15,000 fits satisfy their own discrepancy threshold,
yet null rejection reaches 52.70\%, 96.02\%, and 71.94\%
(Table~\ref{main:tab-inference-size}).  Meeting an estimation threshold
clearly does not establish adequate stopping for inference in these designs.

Figure~\ref{main:fig-inference-stopping} follows the complete adjustment as
the dimension increases.  The top row measures the residual on the sampling
scale; the middle row measures its actual effect on the statistic; and the
bottom row reports the resulting decisions.  These quantities answer
different questions.  A small residual provides a theoretical sufficient
condition, whereas agreement with the direct statistic checks the practical
approximation in the simulated samples.  A rejection rate alone cannot
show whether that approximation is accurate.

\begin{figure}[p]
 \centering
 \includegraphics[width=\textwidth]{stopping_diagnostics.pdf}
 \caption[Effect of stopping on inference]{Stopping diagnostics at $n=100$,
 based on 5,000 samples per model in the separate stopping study.  Rows show
 median $\sqrt n\|e_m\|$, median $|\mathcal T_{n,m}-\mathcal S_n|$, and
 oracle-calibrated null rejection at 5\%.  Circles denote fixed dimensions,
 stars the adaptive variant with the 20-component variance pilot, and
 diamonds the maximum-rank fit.  The dotted line in the top row is the
 median estimation threshold, not a vanishing inferential tolerance.}
 \label{main:fig-inference-stopping}
\end{figure}

\begin{table}[htbp]
 \centering
 \small
 \caption[Stopping and null rejection]{Null rejection percentages at the
 oracle 5\% level.  The adaptive and fixed-$m$ columns use the separate
 stopping experiment; the last column uses the principal null-size samples.
 Each experiment contains 5,000 samples per model at $n=100$.}
 \label{main:tab-inference-size}
 \begin{tabular}{@{}lrrrrr@{}}
 \toprule
 & \multicolumn{4}{c}{Separate stopping experiment}
 & \shortstack{Principal\\size study}\\
 \cmidrule(lr){2-5}\cmidrule(l){6-6}
 Model & \shortstack{Adaptive\\median $m$}
 & \shortstack{Adaptive\\rejection}
 & $m=10$ & $m=70$ & $m=70$\\
 \midrule
 1 & 1 & 52.70 & 6.16 & 6.20 & 5.90\\
 2 & 2 & 96.02 & 5.88 & 6.00 & 5.84\\
 3 & 1 & 71.94 & 5.20 & 5.26 & 6.10\\
 \bottomrule
 \end{tabular}
\end{table}

The adaptive median scaled residuals are 4.2452, 9.4103, and 9.0848.
At $m=70$, they are approximately 0.0031 in every model, and median
absolute differences between $\mathcal T_{n,70}$ and $\mathcal S_n$ are
of order $10^{-5}$.  Continuing to maximum rank changes no rejection
decision.  At $m=10$, median statistic differences are already below 0.09,
although this does not establish ten iterations as a general rule.
Rejection need not approach 5\% monotonically: an intermediate dimension
can give apparently correct size despite a substantial residual.
The residual and statistic difference therefore provide essential checks
alongside rejection frequencies.

For example, stopping at $m=2$ gives 2.64\% rejection in Model~1 but
99.82\% in Model~2.  At $m=5$, all three rejection rates lie between
4.40\% and 5.64\%, although median statistic differences still range from
0.458 to 1.203.  Apparent agreement with the nominal level at one dimension
is consequently weaker evidence than sustained reduction of both the
residual and statistic difference.

In the principal null study, $\mathcal T_{n,70}$ and $\mathcal S_n$ give
identical decisions at the 1\%, 5\%, and 10\% levels in all 15,000 samples.
The 5\% rejection rates of 5.90\%, 5.84\%, and 6.10\% indicate modest
over-rejection; their conditional Monte Carlo standard errors are about
0.33--0.34 percentage points.  The empirical 95th percentiles exceed the
oracle reference values by 3.47\%, 3.87\%, and 3.71\%.  Additional
iterations cannot correct these calibration differences once the fitted
and direct statistics agree.  Models~1 and 2 share the same theoretical
null law because their predictor and error distributions coincide and the
slope cancels under the null.  Their small reported differences arise from
simulation variation.  All reported Monte Carlo standard errors treat the
estimated critical value as fixed and omit its additional uncertainty.

\section{Power against slope changes}
\label{main:sec-inference-power}

For a fixed alternative satisfying $K(\beta-b)\neq0$, moment convergence
and sufficiently late stopping imply
$n^{-1}\mathcal T_{n,m_n}(b)\to_{\mathbb P}\|K(\beta-b)\|^2>0$,
so power tends to one.  Under a shrinking alternative
$\beta_n=b+n^{-1/2}\Delta$, the limit becomes $\|G+K\Delta\|^2$.
Detectability therefore depends on the covariance-weighted departure:
under homoskedasticity,
$\|K\Delta\|^2=\sum_j\lambda_j^2\langle\Delta,v_j\rangle^2$.
These are direction-specific results, not uniform power guarantees over
all slope changes.

In particular, a large unweighted change in the slope can be difficult to
detect when it lies in directions with small covariance eigenvalues.
Conversely, a more modest change in a strongly observed direction can
produce a clear moment shift.  The covariance of the random score also
matters: under general heteroskedasticity, power depends on the orientation
of $K\Delta$ relative to $V$, as well as its norm.  Thus the comparison
between simulation models must retain the particular perturbation used.

The experiment generates
$\beta_\delta(s)=\beta(s)+\delta s$ for
$\delta\in\{-1,-0.95,\ldots,1\}$, and tests the baseline $b=\beta$.
It uses the direct statistic $\mathcal S_n$, evaluated in generating
score coordinates, and a separate empirical 5\% null critical value for
each model and $n\in\{100,200\}$.  These are fixed-perturbation power
curves, not power after 70 CG iterations or a direct experiment with
$\delta=c/\sqrt n$.

\begin{figure}[htbp]
 \centering
 \includegraphics[width=\textwidth]{power_curves.pdf}
 \caption[Direct-moment power]{Power against
 $\beta_\delta(s)=\beta(s)+\delta s$ using $\mathcal S_n$ and empirical-null
 calibration.  Solid curves use $n=100$, dashed curves $n=200$; the dotted
 line marks 5\%.  Each point uses 5,000 separately generated alternative
 samples.}
 \label{main:fig-inference-power}
\end{figure}

At $\delta=0$, rejection lies between 4.58\% and 5.24\% across the six
model--sample-size combinations.  The approximately symmetric curves in
Figure~\ref{main:fig-inference-power} show substantial gains with sample
size.  At $\delta=0.25$, power increases from about 37\% to 65\% in
Models~1 and 2, and from 23.60\% to 44.56\% in Model~3.  At
$\delta=0.50$, the corresponding rates rise from about 90\% to above
99.5\%, and from 75.58\% to 97.26\%.  Models~1 and 2 share covariance
and perturbation direction; Model~3's lower power concerns this particular
direction and does not establish a general ranking.

\section{Confidence sets and their interpretation}
\label{main:sec-inference-confidence}

Inverting a valid level-$\alpha$ test gives
$\{b:\mathcal T_{n,m}(b)\leq q_{1-\alpha}\}$, an asymptotic confidence set
on the identified parameter space.  Unrestricted finite-sample inversion
is unbounded along $\operatorname{Null}(\widehat K)$.  The illustration
therefore restricts candidates to the first five cosine functions,
$b_c(s)=\sum_{j=1}^5c_jh_j(s)$, and uses $m_C=10$ with oracle 5\%
calibration.  If $H$ contains their values on the grid, set
$A=\widehat KH$ and $f=\widehat K\widehat\beta_{10}^{\mathrm{CG}}$.
The exact accepted coefficient set is
\begin{equation}
 C^{(5)}_{0.95}
 =\left\{c\in\mathbb R^5:\frac nT\|f-Ac\|_2^2\leq q_{0.95}\right\}.
 \label{main:eq-inference-coefficient-set}
\end{equation}
For full-rank $A$, it is the ellipsoid
$(c-\widehat c)^\top A^\top A(c-\widehat c)\leq\rho^2$, where
$\widehat c=(A^\top A)^{-1}A^\top f$ and
$\rho^2=Tq_{0.95}/n-\|f-A\widehat c\|_2^2$, or it is empty when
$\rho^2<0$.  The analytic calculation covers all of $\mathbb R^5$,
without the reference illustration's finite coefficient grid and box
restriction.  Appendix~\ref{app:mathematical-proofs} derives the ellipsoid,
its pointwise envelope, and the rank-deficient cases.

Writing $h(s)=(h_1(s),\ldots,h_5(s))^\top$, the exact lower and upper
function values attained by the non-empty ellipsoid are
\begin{equation}
 h(s)^\top\widehat c
 \mathbin{\pm}\sqrt{\rho^2h(s)^\top(A^\top A)^{-1}h(s)}.
 \label{main:eq-inference-envelope}
\end{equation}
These boundaries describe the range of accepted values at each location.
The coefficient vector attaining an extreme can change with $s$, so joining
the extremes need not produce one accepted candidate function.  Membership
is decided by the joint coefficient inequality in
\eqref{main:eq-inference-coefficient-set}.  The envelope is a useful display
of that set, but does not replace its joint constraint.  Similarly, taking
median boundaries across simulated samples describes a typical envelope;
it does not preserve the interpretation of one sample's confidence set.

\begin{table}[htbp]
 \centering
 \small
 \caption[Five-dimensional inversion]{Oracle-calibrated inversion using
 $m_C=10$, $p=5$, and 5,000 samples per model.  Width is averaged over
 the functional grid and then across samples.}
 \label{main:tab-inference-confidence}
 \begin{tabular}{@{}lrrr@{}}
 \toprule
 Model & \shortstack{Full-slope\\non-rejection (\%)}
 & \shortstack{Five-term\\acceptance (\%)}
 & Mean width\\
 \midrule
 1 & 94.20 & 94.32 & 2.810\\
 2 & 94.20 & 94.12 & 2.824\\
 3 & 94.02 & 93.94 & 1.118\\
 \bottomrule
 \end{tabular}
\end{table}

All 15,000 sets are non-empty, bounded, and rank five.
Table~\ref{main:tab-inference-confidence} separates non-rejection of the
full true slope from acceptance of its five-term projection.  The full
truth has non-zero coefficients beyond the fifth term and is outside the
displayed candidate space.  The projected acceptance rates do not establish
coverage for a projection model with an unknown remainder.  Full-slope
acceptance uses the principal null samples with $m=10$ rather than 70;
only 5, 6, and 6 decisions change.  This paired comparison is not independent
coverage evidence.  Model~3's narrower sets concern the selected five
directions, not uniformly smaller uncertainty for the whole slope.

The supplementary figures display pointwise median boundaries across
5,000 envelopes.  Such an aggregate is not a realised 95\% simultaneous
confidence band, and pointwise containment does not establish membership
in the coefficient ellipsoid.  Interpretation also remains conditional on
adequately small iteration error and appropriate calibration.

Together, the results distinguish inadequate stopping from calibration
error and show how sample size improves power along a specified direction.
They concern dense Gaussian curves, homoskedastic errors, and oracle or
simulation-based critical values.  Feasible plug-in and bootstrap
calibration have not been validated here.  Full methodology and results
are retained in Appendices~\ref{chap:inference-methodology}
and~\ref{chap:inference-results}; Chapter~\ref{chap:main-empirical}
therefore assesses prediction on observed spectra without claiming formal
slope inference.
